\documentclass[11pt]{article}
\usepackage[a4paper,margin=18mm]{geometry}
\usepackage{fontspec}
\setmainfont{Latin Modern Roman}
\usepackage{amsmath,amssymb,amsthm,amscd,graphicx,color}
\usepackage{xurl}
\usepackage[unicode,hidelinks]{hyperref}
\allowdisplaybreaks
\setlength\emergencystretch{3em}

\makeatletter
\newtheorem{theorem}{Theorem}
\newtheorem*{theorem*}{Theorem}
\newtheorem{lemma}{Lemma}
\newtheorem*{lemma*}{Lemma}
\newtheorem{corollary}{Corollary}
\newtheorem*{corollary*}{Corollary}
\newtheorem{proposition}{Proposition}
\newtheorem*{proposition*}{Proposition}
\newtheorem{conjecture}{Conjecture}
\newtheorem*{conjecture*}{Conjecture}
{\theoremstyle{definition} \newtheorem{definition}{Definition}
\newtheorem*{definition*}{Definition}
\newtheorem{example}{Example}
\newtheorem*{example*}{Example}
\newtheorem{remark}{Remark}
\newtheorem*{remark*}{Remark}
\newtheorem{note}{Note}
\newtheorem*{note*}{Note}
}

\makeatother
\usepackage{pictexwd,dcpic}

\DeclareMathOperator{\res}{res} \DeclareMathOperator{\Res}{Res}
\DeclareMathOperator{\TR}{TR} \DeclareMathOperator{\Op}{Op}
\DeclareMathOperator{\Tr}{Tr} \DeclareMathOperator{\ord}{ord}
\DeclareMathOperator{\odd}{odd} \DeclareMathOperator{\supp}{supp}
\DeclareMathOperator{\tr}{tr} \DeclareMathOperator{\Log}{Log}
\DeclareMathOperator{\Lie}{Lie} \DeclareMathOperator{\Det}{Det}
\DeclareMathOperator{\Exp}{Exp} \DeclareMathOperator{\rk}{rk}

\def\dbar{d{\hskip-1pt\bar{}}\hskip1pt}
\def\cutoffint{-\hskip -10pt\int}

\newcommand{\R}{{\mathbb R}}
\newcommand{\C}{{\mathbb C}}
\newcommand{\N}{{\mathbb N}}
\newcommand{\Z}{{\mathbb Z}}

\def \Cl {{C\ell}}

\begin{document}
\begin{center}\Large On a closed connected smooth manifold of odd dimension at least three, every complex-linear trace on the full scalar odd-class classical pseudodifferential operator algebra is a scalar multiple of the Kontsevich-Vishik canonical trace\end{center}
\noindent\textbf{Stable ID:} 2140\par
\noindent\textbf{Authors:} Carolina Neira Jimenez; Marie Francoise Ouedraogo\par
\noindent\textbf{Original work:} Classification of Traces and Associated Determinants on Odd-Class Operators in Odd Dimensions\par
\noindent\textbf{Version:} Official SIGMA 8 (2012) article 023, DOI 10.3842/SIGMA.2012.023; journal PDF and native TeX acquired 2026-10-01, exact bytes pinned by SHA256\par
\noindent\textbf{Selected task scope:} Full Theorem2 with standing Section3 hypotheses: scalar operators acting on C-infinity(M), M smooth closed connected, n odd and n>1, hence n>=\allowbreak{}3. Cl\_odd(M) is the union of all integer-order odd-class classical pseudodifferential operators, including smoothing operators and arbitrary positive integer orders. A trace is any complex-linear functional tau vanishing on all commutators in that algebra. Then tau=\allowbreak{}c TR for a complex scalar c, where TR is the Kontsevich-Vishik canonical trace defined by finite-part symbol integration. Fixed nonpositive-order algebras have additional generalized leading-symbol traces and are not this statement.\par
\noindent\textbf{Difficulty:} H2: The author handles the exceptional homogeneous symbol degree minus the dimension by a sphere-Laplacian construction preserving odd-class parity, then globalizes the resulting derivative decomposition to commutators modulo smoothing operators. This analytic-symbol argument, combined with the finite commutator structure of the smoothing ideal, forces every trace to be a multiple of the canonical trace and is substantially beyond ordinary qualification operator theory.\par
\noindent\textbf{Editorial boundary note (not author text):} The full original scalar odd-class canonical-trace uniqueness Theorem 2 is supplied on the entire union of integer-order operators over a closed connected odd-dimensional manifold with n>=\allowbreak{}3. The complete homogeneous-symbol derivative decomposition, including exceptional degree minus n spherical Laplacian/parity and asymptotic summation, localization/partition commutator globalization and rank-one smoothing reduction are retained. No continuity hypothesis is imposed on the selected complex-linear trace. Kontsevich-Vishik canonical-trace existence/invariance/trace properties and the Guillemin smoothing-ideal commutator theorem remain stated cited prior background. Fixed-nonpositive-order classification, its unused positive differential-operator Lemma 5 and all determinant constructions are excluded from the selected task.\par
\noindent\textbf{Source:} \url{https://sigma-journal.com/2012/023/}\quad DOI: \url{https://doi.org/10.3842/SIGMA.2012.023}\par
\noindent\textbf{License and attribution:} Copyright retained by the named authors. Published by SIGMA under CC BY 4.0: \url{https://creativecommons.org/licenses/by/4.0/}. Source: \url{https://sigma-journal.com/about.html}. Adaptation consists of standalone excerpt formatting, cover and numbering restoration; original author language and mathematical text are retained.\par
\medskip\hrule\medskip\noindent\textbf{Author text begins below.}\par
\setcounter{section}{1}
\setcounter{subsection}{0}
\setcounter{subsubsection}{0}
\setcounter{equation}{1}
% Original source lines 87--536
\section{Preliminaries on pseudodif\/ferential operators}\label{section2}

Here we recall the basic notions of classical pseudodif\/ferential operators following \cite{Shubin}.

\subsection{Symbols}

Let $U$ be an open subset of $\mathbb{R}^n$. Given $a\in\mathbb{C}$, a {\it symbol} of order $a$ on $U$ is a complex valued function $\sigma(x,\xi)$ in $C^{\infty}(U\times\mathbb{R}^n)$ such that for any compact subset $K$ of $U$ and any two multiindices $\alpha=(\alpha_1,\dots,\alpha_n), \beta=(\beta_1,\dots,\beta_n) \in \mathbb{N}^n$ there exists a constant $C_{K, \alpha, \beta}$ satisfying for all $(x,\xi)\in K\times\mathbb{R}^n$,
\[
\big| \partial_x^{\alpha}\partial_{\xi}^{\beta}\sigma(x,\xi) \big| \leq C_{K, \alpha, \beta}(1+ \vert \xi \vert )^{\Re(a)- \vert \beta\vert},
\]
where $\partial_x^{\alpha}=\partial_{x_1}^{\alpha_1}\cdots\partial_{x_n}^{\alpha_n}$, $\vert \beta\vert = \beta_1+\cdots+\beta_n$, and $\Re(a)$ stands for the real part of $a$. Let $S^{a}(U)$ denote the set of such symbols.

Notice that if $\Re(a_1)<\Re(a_2)$, then $S^{a_1}(U)\subset S^{a_2}(U)$. We denote by $S^{-\infty}(U):= \bigcap\limits_{a \in \mathbb{C}}S^a(U)$ the space of smoothing symbols on $U$. Given $\sigma\in S^{m_0}(U),\ \sigma_j\in S^{m_j}(U)$, where $m_j\to-\infty$ as $j\to\infty$, we write
\[
\sigma\sim\sum_{j=0}^\infty\sigma_j,
\]
if for every $N\in\mathbb{N}$
\[
\sigma-\sum_{j=0}^{N-1}\sigma_j\in S^{m_N}(U).
\]

The product $\star$ on symbols is def\/ined as follows: if $\sigma_1 \in S^{a_1}(U)$ and $\sigma_2\in S^{a_2}(U)$,
\begin{gather}
\sigma_1\star \sigma_2(x, \xi) \sim \sum_{|\alpha|\geq0}\frac{(-i)^{\vert\alpha\vert}}{\alpha!}\partial_{\xi}^{\alpha} \sigma_1(x,\xi)\partial_x^{\alpha}\sigma_2(x,\xi). \label{eq:starproduct}
\end{gather}
In particular, $\sigma_1\star \sigma_2 \in S^{a_1 + a_2}(U)$.

\subsubsection{Classical symbols}\label{sectionsymbols}

A symbol $ \sigma\in S^a(U)$ is called  {\it classical}, and we write $\sigma\in CS^a(U)$, if there is an asymptotic expansion
\begin{gather}
\sigma(x,\xi) \sim \sum_{j=0}^{\infty}\psi(\xi)\, \sigma_{a-j}(x,\xi). \label{eq:classicallocallogsymb}
\end{gather}
Here $\sigma_{a-j}(x,\xi)$ is a positively homogeneous function in $\xi$ of degree $a-j$:
\[
\sigma_{a-j}(x,t\xi)=t^{a-j}\sigma_{a-j}(x,\xi) \qquad \text{for all}  \ \ t\in\mathbb{R}^+, \ \ \vert \xi\vert \not=0,
\]
and $\psi\in C^{\infty}(\mathbb{R}^n)$ is any cut-of\/f function which vanishes for $\vert \xi \vert \leq \frac{1}{2}$ and such that $\psi(\xi)=1$ for $\vert \xi\vert \geq 1$.

We denote by
\[
CS(U)=\left\langle \bigcup_{a\in \mathbb{C}} CS^a(U)\right\rangle
\]
the algebra generated by all classical symbols on $U$ for the product $\star$.

\subsubsection{Odd-class symbols}

The homogeneous components in the asymptotic expansion of a classical symbol may satisfy some other symmetry relations additional to the positive homogeneity on the second variable. Now we recall the def\/inition of odd-class symbols introduced f\/irst in \cite{kontsevichdeterminant} (see also \cite{grubb}).

\begin{definition}[see \cite{kontsevichdeterminant}]
A classical symbol $\sigma\in {CS}^{a}(U)$ with integer order $a$ is \emph{odd-class} if for each $j\geq 0$, the term $\sigma_{a-j}$ in the asymptotic expansion (\ref{eq:classicallocallogsymb}) satisf\/ies
\begin{gather}
\sigma_{a-j}(x,-\xi)=(-1)^{a-j}\sigma_{a-j}(x,\xi) \qquad {\rm for}\quad \vert \xi\vert \geq 1. \label{eq:OddlogPolyhom}
\end{gather}
In other words, odd-class symbols have the parity one would expect from dif\/ferential operators.
\end{definition}
Let us denote by $CS_{\odd}^a(U)$ the set of odd-class symbols of order $a\in \mathbb{Z}$ on $U$. We set
\[
CS_{\odd}(U)=\bigcup_{a\in \mathbb{Z}} CS_{\odd}^a(U).
\]

\begin{lemma}[see \cite{ducourtiouxthese, kontsevichdeterminant}]\label{lemmaoddclasssymbols}
The odd-class symbols satisfy the following:
\begin{enumerate}\itemsep=0pt
\item The product $\star$ of two odd-class symbols is an odd-class symbol, therefore $CS_{\odd}(U)$ is an algebra.
\item If an odd-class symbol is invertible with respect to the product $\star$, then its inverse is an odd-class symbol.
\end{enumerate}
\end{lemma}

\subsubsection{The noncommutative residue on symbols}

As before, we consider $U$ an open subset of $\R^n$.

\begin{definition}[see \cite{guillemin,wodzicki}]
The \emph{noncommutative residue} of a classical symbol $\sigma\sim\sum\limits_{j=0}^\infty\sigma_{a-j}\in CS^a(U)$ is def\/ined by
\[
\res(\sigma):=\int_U \int_{S_x^{*}U} \sigma_{-n}(x,\xi) \mu(\xi)\, dx=:\int_U \res_x(\sigma)dx,
\]
where $\mu$ is the surface measure on the unit sphere $S_x^{*}U$ over $x$ in the cotangent bundle $T^*U$.
\end{definition}
The noncommutative residue clearly vanishes on symbols of order strictly less than $-n$ and also on symbols of non-integer order.

\begin{lemma}\label{eq: nullityRES}
In odd dimensions, the noncommutative residue of any odd-class symbol vanishes.
\end{lemma}

\begin{proof}
Let $\sigma\in CS_{\odd}^a(U)$ be with asymptotic expansion $\sigma\sim\sum\limits_{j=0}^\infty\psi \sigma_{a-j}$ as in \eqref{eq:classicallocallogsymb}. Since $n$ is odd, we have $\sigma_{-n}(x,-\xi)=(-1)^{n}\sigma_{-n}(x,\xi)=-\sigma_{-n}(x,\xi)$. Then we obtain for any $x\in U$
\begin{gather*}
\res_x(\sigma)=\int_{S_x^{*}U}\!\!\sigma_{-n}(x,\xi)\mu(\xi)=-\int_{S_x^{*}U}\!\!\sigma_{-n}(x,-\xi)\mu(\xi)
=-\int_{S_x^{*}U}\!\!\sigma_{-n}(x,\xi)\mu(\xi)=-\res_x(\sigma).
\end{gather*}
Therefore $\res(\sigma)=0$.
\end{proof}

\subsubsection{An odd-class symbol as a sum of derivatives}

\begin{proposition}[see Lemma~1.3 in \cite{fedosov}, and~\cite{maniccia}]\label{prop:SymbolSumDerivatives}
Let $n\in\Z$ be odd. For any $\sigma\in CS_{\odd}^a(U)$, there exist $\tau_i$ in $CS_{\odd}^{a+1}(U)$ such that
\[
\sigma\sim\sum\limits_{i=1}^n\partial_{\xi_i}\tau_i.
\]
\end{proposition}

\begin{proof}
For a cut-of\/f function $\psi$ as in Section \ref{sectionsymbols} consider
\[
\sigma\sim\sum_{j=0}^\infty\psi \sigma_{a-j},
\]
with $\sigma_{a-j}$ a positively homogeneous function of degree $a-j$ in $\xi$ which satisf\/ies \eqref{eq:OddlogPolyhom}.

 If $a-j\not=-n$, consider the homogeneous function $\tau_{i,a-j+1}:=\frac{\xi_i\sigma_{a-j}(x,\xi)}{a-j+n}$. By Euler's identity we have
\[
\sum_{i=1}^n\partial_{\xi_i}(\tau_{i,a-j+1})(x,\xi)=\sigma_{a-j}(x,\xi).
\]
The homogeneous functions $\tau_{i,a-j+1}$ clearly satisfy \eqref{eq:OddlogPolyhom} for $|\xi|\geq1$:
\begin{gather*}
\tau_{i,a-j+1}(x,t\xi) =t^{a-j+1}\tau_{i,a-j+1}(x,\xi),\qquad \forall\, t>0,\\
\tau_{i,a-j+1}(x,-\xi) =(-1)^{a-j+1}\tau_{i,a-j+1}(x,\xi).
\end{gather*}

 Let $a-j=-n$. In polar coordinates $(r,\omega)\in\mathbb{R}^+\times S^{n-1}$, the Laplacian in $\xi$ reads
\[
\Delta=-\sum_{i=1}^n\partial_{\xi_i}^2=-r^{1-n}\partial_r\big(r^{n-1}\partial_r\big)-r^{-2}\Delta_{S^{n-1}}.
\]
Therefore, for any function $f\in C^\infty(S^{n-1})$,
\[
\Delta\big(f(\omega)r^{2-n}\big)=r^{-n}\Delta_{S^{n-1}}f(\omega).
\]
Since $n$ is odd and $\sigma\in CS_{\odd}^a(U)$, by Lemma \ref{eq: nullityRES} we have $\res(\sigma)=0$. Therefore $\sigma_{-n}(x,\cdot)\upharpoonright_{S^{n-1}}$ is orthogonal to the constants which form the kernel $\ker(\Delta_{S^{n-1}})$. Hence there exists a unique function $h(x,\cdot)\in C^\infty(S^{n-1})$, orthogonal to the constants, such that $\Delta_{S^{n-1}}(h(x,\cdot))=\sigma_{-n}(x,\cdot)\upharpoonright_{S^{n-1}}$. The function $h(x,-\xi)+h(x,\xi)$ is constant and orthogonal to the constants, therefore, $h(x,\cdot)$ is an odd function on $S^{n-1}$.

 We choose a smooth function $\chi$ on $\mathbb{R}$ which vanishes for small $r$ and is equal to 1 for $r\geq1/2$. For $r=|\xi|$, we set
\[
b_{-n}(x,\xi):=\chi(|\xi|)|\xi|^{2-n}h\left(x,\frac{\xi}{|\xi|}\right).
\]
The function $b_{-n}$ is smooth on $U\times\mathbb{R}^n$ and is homogeneous of degree $-n+2$ in $\xi$ for $|\xi|\geq1$. As $\sigma_{-n}(x,\xi)$ vanishes for $x$ outside a compact set, so does $b_{-n}(x,\xi)$. In particular, $b_{-n}$ is a symbol of order $-n+2$ on $U$. Let us def\/ine $\tau_{i,-n+1}:=-\partial_{\xi_i}b_{-n}$. Since $h$ is odd so is $b_{-n}$ and therefore,
\[
\tau_{i,-n+1}(x,-\xi)=-(\partial_{\xi_i}b_{-n})(x,-\xi)
=-\partial_{\xi_i}b_{-n}(x,\xi)
=(-1)^{-n+1}\tau_{i,-n+1}(x,\xi).
\]
Moreover, we have for $r=|\xi|\geq1$
\[
\Delta b_{-n}(x,\cdot)=\Delta\big(r^{2-n}h(x,\cdot)\big)=r^{-n}\left(\sigma_{-n}(x,\cdot)\upharpoonright_{S^{n-1}}\right)=\sigma_{-n}(x,\cdot).
\]


Let $\tau_i\sim\sum\limits_{j=0}^\infty\psi\,\tau_{i,a-j+1}$, then since $\partial_{\xi_i}\psi$ has compact support, the dif\/ference $\sigma-\sum\limits_{i=1}^n\partial_{\xi_i}\tau_i$ is smoothing and
\begin{gather*}
\sigma\sim\sum_{i=1}^n\sum_{j=0}^\infty\psi \partial_{\xi_i}\tau_{i,a-j+1}\sim\sum_{i=1}^n\partial_{\xi_i}\tau_i. \tag*{\qed}
\end{gather*}
\renewcommand{\qed}{}
\end{proof}

\subsection{Pseudodif\/ferential operators}

Let $U\subset \mathbb{R}^n$ be an open subset, and denote by $C^{\infty}_c (U)$ the space of smooth compactly supported functions on $U$. To the symbol $\sigma\in S(U)$, we associate the linear integral operator $\Op(\sigma):C^{\infty}_c (U)\to C^{\infty}(U)$ def\/ined for $u\in C^{\infty}_c (U)$ by
\begin{gather*}
\Op(\sigma)(u)(x) =\int_{T_x^*U}{e^{ix\cdot\xi}\sigma(x, \xi)\widehat u(\xi)\dbar\xi}
 =\int_{T_x^*U}{\int_{U}{e^{i(x-y)\cdot\xi}\sigma(x,\xi) u(y) dy\dbar\xi}}\\
\hphantom{\Op(\sigma)(u)(x)}{} =\int_{U}k(x,y)u(y)dy,
\end{gather*}
where $\widehat u(\xi)=\int_{U}{e^{-iy\cdot\xi}u(y)\,dy}$ is the Fourier transform of $u$ and $\dbar\xi:=(2\pi)^{-n}d\xi$.
In this expression $k(x,y)=\int{e^{i(x-y)\cdot\xi}\sigma(x, \xi)\dbar\xi}$ is seen as a distribution on $U\times U$ that is smooth outside the diagonal. We say that $\Op(\sigma)$ is a {\it pseudodifferential operator} ($\psi$DO) with Schwartz kernel given by $k(x,y)$. An operator is \emph{smoothing} if its Schwartz kernel is a smooth function on $U\times U$. If $\sigma\sim \sum\limits_{j=0}^{\infty}\psi\sigma_{a-j}$ is a classical symbol of order~$a$, then $A=\Op(\sigma)$ is called a {\it classical pseudodifferential operator} of order~$a$. The homogeneous component $\sigma_a$ of $\sigma$ is called {\it the leading symbol} of~$A$, and will be denoted by $\sigma^L_A$.

A $\psi$DO $A$ on $U$ is called {\it properly supported} if for any compact $K\subset U$, the set $\{(x,y)\in \supp(k_A): x\in K  \ {\rm or}\ y\in K\}$ is compact, where $\supp(k_A)$ denotes the support of the Schwartz kernel of~$A$.  Any $\psi$DO $A$ can be written in the form (see~\cite{Shubin})
\begin{gather}
A=P+R, \label{eq:PDO}
\end{gather}
where $P$ is properly supported and $R$ is a smoothing operator.

A properly supported $\psi$DO maps $C^{\infty}_c(U)$ into itself. The product $\star$ on symbols def\/ined in \eqref{eq:starproduct} induces a composition on properly supported $\psi$DOs on $U$. The composition $AB$ of two properly supported $\psi$DOs $A$ and $B$ is a properly supported $\psi$DO with symbol $\sigma(AB)=\sigma(A) \star \sigma(B)$.

The notion of a $\psi$DO can be extended to operators acting on manifolds (see Section~4.3 in~\cite{Shubin}). Let $M$ be a smooth closed manifold of dimension $n$. A linear operator $A:C^{\infty}(M)\to C^{\infty}(M)$ is a \emph{pseudodifferential operator} of order $a$ on $M$ if in any atlas, $A$ is locally a pseudodif\/ferential operator. This means that given a local coordinate chart $U$ of $M$, with dif\/feo\-mor\-phism $\varphi:U\to V$, from $U$ to an open set $V\subseteq\R^n$, the operator $\varphi^\#A$ def\/ined by the following diagram is a pseudodif\/ferential operator of order $a$ on $V$:
\[
\begindc{\commdiag}[35]
\obj(1,3)[objLie(cal{G})]{$C_c^\infty(V)$}
\obj(4,3)[objMathbb{C}]{$C^\infty(V)$}
\obj(4,1)[objMathbb{C}^*]{$C^\infty(U)$}
\obj(1,1)[objCal{G}]{$C_c^\infty(U)$}
\mor{objLie(cal{G})}{objMathbb{C}}{$\varphi^\#A$}
\mor{objMathbb{C}}{objMathbb{C}^*}{$\varphi^*$}
\mor{objCal{G}}{objMathbb{C}^*}{$r_{U}\circ A\circ i_{U}$}
\mor{objLie(cal{G})}{objCal{G}}{$\varphi^*$} [\atright,\solidarrow]
\enddc
\]
where $i_{U}:C_c^\infty(U)\to C_c^\infty(M)$ is the natural embedding, and $r_{U}:C^\infty(M)\to C^\infty(U)$ is the natural restriction.

Let $\Cl^a(M)$ denote the set of classical $\psi$DOs of order $a$ on $M$, i.e.\ operators whose symbol is classical of order $a$ in any local chart of $M$. If $A_1\in \Cl^{a_1}(M), A_2\in \Cl^{a_2}(M)$, then $A_1 A_2\in \Cl^{a_1 + a_2}(M)$, thus, the space $\Cl^{a}(M)$ is an algebra if and only if $a$ is an integer and $a\leq0$. We denote by
\[
\Cl(M):= \left\langle\bigcup_{a\in \mathbb{C}} \Cl^a(M)\right\rangle
\]
the algebra generated by all classical $\psi$DOs on $M$, and by
\[
\Cl^{-\infty}(M):=\bigcap_{a\in \mathbb{C}} \Cl^a(M)
\]
the ideal of smoothing operators in $\Cl(M)$.

We will also denote by
\[
\Cl^{\notin\Z}(M):=\left\langle\bigcup\limits_{a\notin\Z\cap[-n,+\infty)}\Cl^{a}(M)\right\rangle
\]
the space generated by classical $\psi$DOs on $M$ whose order is non-integer or less than $-n$.

A classical operator $A\in \Cl^{a}(M)$ of integer order $a$ is odd-class if in any local chart its local symbol $\sigma(A)$ is odd-class.
We denote by $\Cl_{\odd}^a(M)$ the set of odd-class operators of order $a\in \mathbb{Z}$ and we def\/ine
\[
\Cl_{\odd}(M)=\bigcup_{a\in \mathbb{Z}} \Cl_{\odd}^a(M).
\]

As in Lemma \ref{lemmaoddclasssymbols}, the following lemma implies that $\Cl_{\odd}(M)$ is an algebra:
\begin{lemma}[see Section~4 in~\cite{kontsevichdeterminant}]\label{lema:oddClassproduct}
Let $A\in \Cl^{a}_{\odd}(M)$ and $B\in \Cl^{b}_{\odd}(M)$, $a,b\in\Z$. Then $AB\in \Cl^{a+b}_{\odd}(M)$. If moreover $B$ is an invertible elliptic operator, then $B^{-1}\in \Cl^{-b}_{\odd}(M)$ and $AB^{-1}\in \Cl^{a-b}_{\odd}(M)$.
\end{lemma}

The algebra $\Cl_{\odd}(M)$ contains the dif\/ferential operators and their parametrices.

\begin{remark}\label{remarknotationClodd}
Even though the def\/inition of odd-class pseudodif\/ferential operators makes sense on any closed manifold, in this paper we restrict ourselves to odd-dimensional closed manifolds. The reason is that the canonical trace (which will be explained below in Section~\ref{subsectioncanonicaltrace}) is well def\/ined only in that case. So, from now on, the notation $\Cl_{\odd}(M)$ will be used only when the dimension $n$ of the manifold $M$ is odd.
\end{remark}

\subsubsection{Fr\'echet topology on pseudodif\/ferential operators}\label{subsectionfrechettop}

For any complex number $a$,  we equip the vector space $\Cl^a(M)$ with a Fr\'echet topology as follows. Let us consider a covering of $M$ by open neighborhoods $\{U_i\}_{i\in I}$, a f\/inite subordinated partition of unity $(\chi_i)_{i\in I}$ and smooth functions $(\widetilde{\chi_i})_{i\in I}$ on $M$ such that $\supp(\widetilde{\chi_i})\subset U_i$ and $\widetilde{\chi_i}=1$ near the support of $\chi_i$. By \eqref{eq:PDO} any $\psi$DO $A$ can be written in the form $A=\sum\limits_{i\in I}(A_i+R_i)$ where the operators $A_i:=\chi_i\cdot\Op(\sigma_i)\cdot\widetilde{\chi_i}\in Cl^{a}(M)$ are properly supported in $U_i$ with symbols $\sigma^{(i)}(A)\sim\sum\limits_{j=0}^\infty\psi \sigma_{a-j}^{(i)}(A)$, and $R_i$ is a smoothing operator with smooth kernel $k_i$ which has compact support in $U_i\times U_i$.

We equip ${\Cl}^{a}(M)$ with the following countable set of semi-norms: for any compact subset $K\subset U_i$ for any $j\geq 0$, $N\geq 1$ and for any multiindices $\alpha$, $\beta$
\begin{gather*}
%\label{eq:SeminormLOGPOLY}
 \sup\limits_{x\in K}\sup\limits_{
  \xi\in\mathbb{R}^n}(1+\vert\xi\vert)^{\vert\beta\vert-a}
  \big|\partial_x^{\alpha}\partial_{\xi}^{\beta} \sigma^{(i)}(A)(x,\xi)\big|,\\
 \sup\limits_{x\in K}\sup\limits_{
  \xi\in\mathbb{R}^n}(1+\vert\xi\vert)^{\vert\beta\vert-a+N}
  \left|\partial_x^{\alpha}\partial_{\xi}^{\beta} \left(\sigma^{(i)}(A)-\sum_{j=0}^{N-1}\chi(\xi)  \sigma_{a-j}^{(i)}(A)\right)(x,\xi)\right|,\\
 \sup\limits_{x\in K}\sup\limits_{ \vert\xi\vert=1}
  \big|\partial_x^{\alpha}\partial_{\xi}^{\beta}\sigma_{a-j}^{(i)}(A)(x,\xi)\big|,\qquad
 \sup\limits_{x,y\in K}
  \big|\partial_x^{\alpha}\partial_{y}^{\beta} k_i(x,y)\big|.
\end{gather*}

\subsubsection{The logarithm of a classical pseudodif\/ferential operator}

An operator $A \in \Cl(M)$ with positive order has \emph{principal angle} $\theta$ if for every $(x, \xi)\in T^*M\setminus\{0\}$, the leading symbol $\sigma_A^L(x, \xi)$ has no eigenvalues on the ray $L_\theta=\{re^{i\theta}, r\geq 0\}$; in that case $A$ is \emph{elliptic}.
\begin{definition}[see e.g.~\cite{okikiolu}]
An operator $A\in \Cl(M)$ is \emph{admissible} with spectral cut $\theta$ if $A$ has principal angle $\theta$ and the spectrum of $A$ does not meet $L_{\theta}=\{re^{i\theta}, r\geq 0\}$. In particular such an operator is invertible and elliptic. The angle $\theta$ is called an \emph{Agmon angle} of $A$.
\end{definition}
Let $A\in \Cl(M)$ be admissible with spectral cut $\theta$ and positive order $a$. For $\Re(z)<0$, the complex power $A_{\theta}^z$ of $A$ is
def\/ined by the Cauchy integral (see \cite{seeley})
\[
A_\theta^{z}=\frac{i}{2\pi} \int_{\Gamma_{r,\theta}} \lambda_\theta^{z} (A-\lambda)^{-1}\, d\lambda,
\]
where $\lambda_\theta^z= \vert \lambda\vert^z e^{iz {(\rm arg}\lambda)}$ with  $\theta\leq {\rm arg}\,\lambda<\theta+2\pi$. Here $\Gamma_{r,\theta}$ is a contour along the ray $L_\theta$ around the (non-zero) spectrum of $A$, and $r$ is any small positive real number such that $\Gamma_{r,\theta}$ does not meet the spectrum of $A$. The operator $A_{\theta}^z$ is an elliptic classical $\psi$DO of order $az$; in particular, for $z=0$, we have $ A_{\theta}^0 = I$.

The def\/inition of complex powers can be extended to the whole complex plane by setting $A_{\theta}^z:=A^kA_{\theta}^{z-k}$ for $k\in \N$ and $\Re(z)<k$; this def\/inition is independent of the choice of $k$ in $\N$ and preserves the usual properties, i.e.\ $A_{\theta}^{z_1}A_{\theta}^{z_2}=A_{\theta}^{z_1+z_2}$, $A_{\theta}^k=A^k$, for $k\in\Z$.
Complex powers of operators depend on the choice of spectral cut. Indeed, if $L_{\theta}$ and $L_{\phi}$ are two spectral cuts for~$A$ outside an angle which contains the spectrum of $\sigma^L(A)(x,\xi)$ then
\begin{gather}
A_\theta^z-A_\phi^z=\big(1-e^{2i\pi z} \big)  P_{\theta, \phi} (A) A_\theta^z, \label{eq:DifferenceComplexPowers}
\end{gather}
where the operator
\begin{gather}
 P_{\theta, \phi} (A)=\frac{1}{2i\pi} \int_{\Gamma_{\theta,\phi}}\lambda^{-1}A(A-\lambda)^{-1}\,d\lambda \label{projection}
\end{gather}
is a projection (see \cite{PongeAsymmetry, WodzickiThesis}). Here $\Gamma_{\theta,\phi}$ is a contour separating the part of the spectrum of $A$ contained in the open sector $\theta<\arg\lambda<\phi$ from the rest of the spectrum.

The logarithm  of an admissible operator $A$ with spectral cut $\theta$  is def\/ined in terms of the derivative at $z=0$ of its complex power:
\[
\log_\theta(A)= \partial_z {A_\theta^{z}} _{\vert_{z=0}}.
\]
Logarithms of classical $\psi$DOs of positive order are not classical anymore since their symbols involve a logarithmic term $\log\vert \xi\vert$ as the following elementary result shows.

\begin{proposition}[\cite{okikiolu}]\label{prop:leadsymbolsigmathetaA}
Let $A\in \Cl^{a}(M, E)$ be an admissible operator with positive order and spectral cut $\theta$. In a local trivialization, the symbol of $\log_\theta(A)$ reads:
\[
\sigma_{\log_\theta(A)}(x, \xi)=a  \log \vert \xi\vert  I +\sigma_0^A(x,\xi),
\]
where $\sigma_0^A$ is a classical symbol of order zero.
\end{proposition}

\begin{remark}
If $A$ is a classical $\psi$DO of order zero then $A$ is bounded, and if it admits a spectral cut, then complex powers and the logarithm of~$A$ are directly def\/ined using a Cauchy integral formula, and they are classical $\psi$DOs (see~\cite{okikiolu} and Remark~2.1.7 in~\cite{ouedraogothese}).
\end{remark}

Just as complex powers, the logarithm depends on the choice of spectral cut. Indeed, given two spectral cuts $\theta,\, \phi$ of the operator $A$ such that $0\leq \theta <\phi <2\pi$, dif\/ferentiation of (\ref{eq:DifferenceComplexPowers}) with respect to $z$ and evaluation at $z=0$ yields
\begin{gather}
\log_\theta A-\log_\phi A=-2i\pi P_{\theta, \phi}(A). \label{eq:difflogarithms}
\end{gather}

\subsubsection{Pseudodif\/ferential operators in terms of commutators}

In this subsection we use the $\psi$DO analysis techniques similar to the ones implemented in \cite{maniccia}; we assume that $M$ is an $n$-dimensional closed manifold and $n$ is odd. Given a function $f\in C^\infty(M)$, we also denote by $f$ the zero-order classical $\psi$DO given by multiplication by $f$.
\begin{proposition}\label{prop:PDOsumComutators}
If $A\in\Cl^{a}_{\odd}(M)$, then there exist functions $\alpha_k\in C^\infty(M)$, operators $B_k$ in $\Cl^{a+1}_{\odd}(M)$ and a smoothing operator $R_A$ such that
\begin{gather}
A=\sum_{k=1}^N[\alpha_k,B_k]+R_A. \label{eq:PDOsumComutators}
\end{gather}
\end{proposition}

\begin{proof}
Let us consider a covering of $M$ by open neighborhoods $\{U_j\}_{j=1}^N$ and a f\/inite subordinated partition of unity $\{\varphi_j\}_{j=1}^N$, such that for every pair $(j,k)$, both $\varphi_j$ and $\varphi_k$ have support in one coordinate neighborhood. We write (see Subsection \ref{subsectionfrechettop})
\[
A=\sum_{j,k}\varphi_jA\varphi_k+R.
\]
Each operator $\varphi_jA\varphi_k$ may be considered as an odd-class $\psi$DO on $\mathbb{R}^n$ with symbol $\sigma$ in $CS^a_{\odd}(\mathbb{R}^n).\!$ By Proposition \ref{prop:SymbolSumDerivatives}, there exist odd-class symbols $\tau_l$ of order $a+1$ such that
\[
\sigma\sim\sum\limits_{l=1}^n\partial_{\xi_l}\tau_l.
\]
For any symbol $\tau$ we have,
\[
\sigma([x_l,\Op(\tau)])\sim x_l\cdot \tau-\tau\cdot x_l-i^{-1}\partial_{\xi_l}\tau =i\partial_{\xi_l}\tau,
\]
so that
\[
\Op(\partial_{\xi_l}\tau)=-i[x_l,\Op(\tau)] \quad {\rm up\ to\ a\ smoothing\ operator.}
\]
Since $\sigma\sim\sum\limits_{l=1}^n\partial_{\xi_l}\tau_l,$ there exist smoothing operators $R''$, $R^\prime$ such that
\[
\Op(\sigma)=\Op\left(\sum\limits_{l=1}^n\partial_{\xi_l}\tau_l\right)+R''=-i\sum\limits_{l=1}^n [x_l,\Op(\tau_l)]+R^\prime.
\]
For each index $j$ let $\psi_j\in C_c^\infty(U_j)$ be such that $\psi_j\equiv1$ near $\supp(\varphi_j)$. Then for $\chi$ in $C_c^\infty(\mathbb{R}^n)$ such that $\chi\varphi_j=\varphi_j$, $\chi\psi_j=\psi_j$, we have
\begin{gather*}
\varphi_j\Op(\sigma)\psi_j
 =-i\sum\limits_{l=1}^n\varphi_j[x_l,\Op(\tau_l)]\psi_j+\varphi_jR^\prime\psi_j
 =-i\sum\limits_{l=1}^n [\chi x_l,\varphi_j\Op(\tau_l)\psi_j]+\varphi_jR^\prime\psi_j.
\end{gather*}
As in (\ref{eq:PDO}) we write $\varphi_j A \psi_j= \Op(\sigma_{j}) +R_{j}$ for some $\sigma_j\in CS^{a}_{\odd}(\mathbb{R}^n)$ and some smoothing ope\-ra\-tor~$R_{j}$. We have $A= \sum _{j}\Op(\sigma_{j})+\sum _{j}R_{j}+R$, hence $A$ can be written in the form
\[
A=\sum\limits_{k=1}^{N}[\alpha_k,B_k]+R_A,
\]
where $\alpha_k$ is a smooth function on $M$ (and represents the operator in $\Cl^{0}_{\odd}(M)$ of multiplication by $\alpha_k$), $B_k$ lies in $\Cl^{a+1}_{\odd}(M),$ and $R_A$ is a smoothing operator.
\end{proof}

\begin{corollary}\label{corol2}
If $A\in \Cl^{a}_{\odd}(M)$ then there exist $B_i\in \Cl^{a}_{\odd}(M)$, smooth functions $a_i\in C^\infty(M)$, and a smoothing operator $R$ such that
\[
A-\Op\big(\sigma_A^L\big)=\sum_{i=1}^n[a_i,B_i]+R.
\]
\end{corollary}

\begin{proof}
It follows by applying Proposition~\ref{prop:PDOsumComutators} to $A-\Op(\sigma_A^L)\in \Cl^{a-1}_{\odd}(M)$.
\end{proof}

\section{Classif\/ication of traces on odd-class operators}\label{sectiontraces}

The classif\/ication made in this section is essentially based on Proposition \ref{prop:PDOsumComutators}, namely the decomposition of an odd-class operator in terms of commutators of odd-class operators. Let us f\/irst recall the def\/inition of a trace.

Let $M$ be a closed connected manifold of dimension $n>1$, and let $\mathcal{A}\subseteq\Cl(M)$. A~\emph{trace} on~$\mathcal{A}$ is a map
\[
\tau: \ \mathcal{A}\rightarrow\mathbb{C},
\]
linear in the sense that for all $a, b\in\C$, whenever $A$, $B$ and $aA+bB$ belong to $\mathcal{A}$ we have
\[
\tau(aA+bB)=a \tau(A)+b \tau(B),
\]
and such that for any $A,B\in\mathcal{A}$, whenever $AB,BA\in\mathcal{A}$ it satisf\/ies
\[
\tau([A,B])=0,\qquad \text{or equivalently,} \qquad \tau(AB)=\tau(BA).
\]

\subsection{Examples of traces on pseudodif\/ferential operators}

Interestingly by Lemma \ref{eq: nullityRES}, the noncommutative residue, which is the only trace on the whole space $\Cl(M)$ (see \cite{fedosov,lesch,wodzicki}), vanishes on odd-class operators when the dimension of the manifold is odd. In this section we review the traces which are non-trivial on this class of operators.

\subsubsection[The $L^2$-trace]{The $\boldsymbol{L^2}$-trace}\label{sectionL2trace}

A $\psi$DO $A$ whose order has real part less than $-n$ is a trace-class operator. The $L^2$-trace (also called operator trace) is the functional
\begin{gather}
\Tr:\ \left\langle\bigcup\limits_{\Re(a)<-n}\Cl^{a}(M)\right\rangle \to\mathbb{C},\qquad
A \mapsto\Tr(A):=\int_{M}k_A(x,x)\,dx, \label{TrL2kernel}
\end{gather}
where $k_A$ is the Schwartz kernel of the operator $A$. This trace is continuous for the Fr\'echet topology on the space of $\psi$DOs of constant order less than $-n$.

This is the unique trace on the algebra of smoothing operators $\Cl^{-\infty}(M)$, since we have the exact sequence (see~\cite{guilleminfourier})
\[
0\to\big[\Cl^{-\infty}(M),\Cl^{-\infty}(M)\big]\to \Cl^{-\infty}(M)\overset{\Tr}{\longrightarrow}\mathbb{C}\to0.
\]
More precisely we have

\begin{theorem}[Theorem~A.1 in~\cite{guilleminfourier}]\label{smooth}
If $R$ is a smoothing operator then, for any pseudodifferential idempotent $J$, of rank~$1$, there exist smoothing operators $S_1,\ldots,S_{N}$, $T_1,\ldots,T_{N}$, such that
\[
R=\Tr(R)J+\sum_{j=1}^{N}[S_j,T_j].
\]
Therefore, any smoothing operator with vanishing $L^2$-trace is a sum of commutators in the space $[\Cl^{-\infty}(M),\Cl^{-\infty}(M)]$.
\end{theorem}

\begin{proposition}[see e.g.\ Introduction in~\cite{lesch} and Proposition~4.4 in~\cite{leschtraces}]\label{nocont}
The trace~$\Tr$ does not extend to a trace functional neither on the whole algebra $\Cl(M)$, nor does it on the algebra~$\Cl^0(M)$.
\end{proposition}

\subsubsection{The canonical trace}\label{subsectioncanonicaltrace}

We start with the def\/inition of the cut-of\/f integral of a symbol, as in \cite{sylvie}. Let $U$ be an open subset of $\mathbb{R}^n$.

\begin{proposition}[see \cite{sylvie} and Section~1.2 in~\cite{scott}]
Let $\sigma\in CS^a(U)$, such that for any $N\in\mathbb{N}$, $\sigma$ can be written in the form $ \sigma(x,\xi) =\sum\limits_{j=0}^{N-1}\psi(\xi)\sigma_{a-j}(x,\xi)+\sigma_N(x,\xi)$, with $\sigma_{a-j}$, $\psi$ and $\sigma_N\in S^{a-N}(U)$ as in~\eqref{eq:classicallocallogsymb}. The integral of $\sigma$ over $B^{*}_x(0,R)$ $($which stands for the ball of radius $R$ in the cotangent space $T^{*}_xU)$ has the asymptotic expansion
\[
\int_{B^{*}_x(0,R)}\sigma(x,\xi)d\xi \underset{R\to\infty}{\sim}\sum_{\substack{j=0\\ a-j+n\not=0}}^{N-1}\alpha_j(\sigma)(x)R^{a-j+n} +\res_x(\sigma)\log R+\alpha_x(\sigma),
\]
where $\alpha_x(\sigma)$ converges when $R\to\infty$.
\end{proposition}

\begin{proof}
We write
\[
\int_{B_x^*(0,R)}\psi(\xi)\sigma_{a-j}(x,\xi)d\xi=\int_{B_x^*(0,1)}\psi(\xi)\sigma_{a-j}(x,\xi)d\xi +\int_{B_x^*(0,R)\setminus B_x^*(0,1)}\psi(\xi)\sigma_{a-j}(x,\xi)d\xi.
\]
Using the fact that $\sigma_{a-j}$ is positively homogeneous of degree $a-j$ we have
\begin{gather*}
 \int_{B_x^*(0,R)\setminus B_x^*(0,1)}\psi(\xi)\sigma_{a-j}(x,\xi)d\xi=\int_1^R\int_{\vert\omega\vert=1}r^{a-j+n-1}\sigma_{a-j}(x,\omega)\,d\omega dr \\
\qquad{} =\dfrac{1}{a-j+n}R^{a-j+n} \int_{\vert\omega\vert=1}\sigma_{a-j}(x,\omega)d\omega-\dfrac{1}{a-j+n}\int_{\vert\omega\vert=1}\sigma_{a-j}(x,\omega)\,d\omega.
\end{gather*}
It follows that the integral $\int_{B_x^*(0,R)}\sigma(x,\xi)\,d\xi$ admits the asymptotic expansion
\begin{gather*}
\int_{B_x^*(0,R)}\sigma(x,\xi)\,d\xi\underset{R\to\infty}{\sim}\sum_{\substack{j=0\\a-j+n\not=0}}^{N-1} \frac{1}{a-j+n}R^{a-j+n}\int_{\vert\omega\vert=1}\sigma_{a-j}(x,\omega)\,d\omega \\
\hphantom{\int_{B_x^*(0,R)}\sigma(x,\xi)\,d\xi\underset{R\to\infty}{\sim}}{}
+{\res}_x(\sigma)\log R + \alpha_x(\sigma),
\end{gather*}
where $\alpha_x(\sigma)$ is given by
\begin{gather}
\alpha_x(\sigma) := \int_{B_x^*(0,R)}\sigma_N(x, \xi)\,d\xi+\sum_{j=0}^{N-1} \int_{B_x^*(0,1)} \psi(\xi)\sigma_{a-j}(x,\xi)\,d\xi\nonumber\\
\hphantom{\alpha_x(\sigma) :=}{}
-\sum_{\substack{j=0\\ a-j+n\not=0}}^{N-1} \dfrac{1}{a-j+n}\int_{\vert\omega\vert=1}\sigma_{a-j}(x,\omega)\,d\omega,\label{eq:CutoffIntegral}
\end{gather}
which may depend on the variable $x$.
\end{proof}

 For $N$ suf\/f\/iciently large, the integral $\int_{B_x^*(0,R)}\sigma_N(x,\xi)d\xi$ is well def\/ined, so the term $\alpha_x(\sigma)$ given by~\eqref{eq:CutoffIntegral} converges when $R\to\infty$, and taking this limit we def\/ine the \emph{cut-off integral} of $\sigma$ by
\begin{gather*}
{-\hskip -11,5pt\int}_{T^{*}_xU}\sigma(x,\xi)\,d\xi :=\int_{T^{*}_xU}\sigma_N(x,\xi)\,d\xi+\sum_{j=0}^{N-1} \int_{B^{*}_x(0,1)}\psi(\xi)\sigma_{a-j}(x,\xi)\,d\xi\\
\hphantom{{-\hskip -11,5pt\int}_{T^{*}_xU}\sigma(x,\xi)\,d\xi :=}{}  -\sum_{\substack{j=0\\ a-j+n\not=0}}^{N-1}\dfrac{1}{a-j+n}\int_{S^{*}_xU}\sigma_{a-j}(x,\omega)\,d\omega,
\end{gather*}
where $S^{*}_xU$ stands for the unit sphere in the cotangent space $T^{*}_xU$. This def\/inition is independent of $N>a+n-1$. Moreover, if $a<-n$, then ${-\hskip -10pt\int}_{\mathbb{R}^n}\sigma(x,\xi)d\xi=\int_{\mathbb{R}^n}\sigma(x,\xi)d\xi$.

According to Proposition \ref{nocont}, there is no non-trivial trace on $\Cl(M)$ which extends the $L^2$-trace. However, the $L^2$-trace does extend to non-integer order operators and to odd-class operators. Indeed, M.~Kontsevich and S.~Vishik~\cite{kontsevichdeterminant} constructed such an extension, the \emph{canonical trace}
\begin{gather*}
\TR: \ \Cl^{\notin\Z}(M)\bigcup\Cl_{\odd}(M) \to\mathbb{C},\qquad
 A  \mapsto\TR(A):=\dfrac{1}{(2\pi)^n}\int_Mdx\ {-\hskip -11,5pt\int}_{T_x^{*}M}\sigma(A)(x,\xi)d\xi,
\end{gather*}
where the right hand side is def\/ined using a f\/inite covering of $M$, a partition of unity subordinated to it and the local representation of the symbol, but this def\/inition is independent of such choices. As we already stated in Remark~\ref{remarknotationClodd}, the canonical trace is well def\/ined on $\Cl_{\odd}(M)$ only when the dimension $n$ of the manifold is odd (see~\cite{grubb,lesch}), which is always our case.

 If $A\in \Cl^{a}(M)$, $B\in \Cl^b(M)$ and if $a,b\notin\Z$, then $\ord(AB)=a+b$ may be an integer, so the linear space $\Cl^{\notin\Z}(M)$ is not an algebra; in spite of this, the canonical trace has the following properties (see~\cite{kontsevichdeterminant}, Section~5 in~\cite{lesch}, and \cite{ouedraogo,sylvie,scott}):
 \begin{enumerate}\itemsep=0pt
\item For any $c\in\mathbb{C}$, if $A,B\in\Cl^{\notin\Z}(M)$ are such that $\ord(c A+B)\notin\Z\cap[-n,+\infty)$, or if $A,B\in\Cl_{\odd}(M)$, then $\TR(c A+B)=c\TR(A)+\TR(B)$.
\item For any $A\in\Cl(M)$ such that $\ord(A)<-n$, $\TR(A)=\Tr(A)$, i.e.\ the canonical trace extends the $L^2$-trace def\/ined in \eqref{TrL2kernel}.
\item If $A,B\in\Cl^{\notin\Z}(M)$ are such that $AB\in\Cl^{\notin\Z}(M)$, or if $A,B\in\Cl_{\odd}(M)$, then $\TR(AB)=\TR(BA)$.
\end{enumerate}

The canonical trace is continuous for the Fr\'echet topology on the sets of $\psi$DOs of constant order in $\Cl^{\notin\Z}(M)\bigcup\Cl_{\odd}(M)$.

\setcounter{section}{3}
\setcounter{subsection}{1}
\setcounter{subsubsection}{3}
\setcounter{equation}{12}
% Original source lines 566--597
\subsection[Trace on $\Cl_{\odd}(M)$]{Trace on $\boldsymbol{\Cl_{\odd}(M)}$}

The canonical trace is the unique trace (up to a constant) on its domain. This result was proved in~\cite{maniccia} (which goes back to 2007 but it was published in 2008) using the fact that the operator corresponding to the derivative of a symbol is, up to a smoothing operator, proportional to the commutator of appropriate operators. The latter idea was considered in~\cite{sylvie} to show the equivalence between Stokes' property for linear forms on symbols and the vanishing of linear forms on commutators of operators. In~\cite{ponge} the author uses the Schwartz kernel representation of an operator to express any non-integer order operator and any operator of regular parity class as a sum of commutators up to a smoothing operator, and then to give another proof of the uniqueness of the canonical trace. The following proof, that we f\/ind in Proposition~3.2.4 of~\cite{ouedraogothese}, is done in the spirit of~\cite{maniccia}.

\begin{theorem}\label{prop:UniquenessTrace}
Any trace on $\Cl_{\odd}(M)$ is proportional to the canonical trace.
\end{theorem}

\begin{proof}
By \eqref{eq:PDOsumComutators} any operator $A$ in $\Cl_{\odd}(M)$ can be written in the form
\[
A=\sum\limits_{k=1}^{N}[\alpha_k,B_k]+R_A,
\]
where $\alpha_k$ are smooth functions on $M$ that can be seen as elements of $\Cl^0_{\odd}(M)$, $B_k$ belong to $\Cl^{a+1}_{\odd}(M),$ and $R_A$ is a smoothing operator.
By Theorem \ref{smooth}, we can express $R_A$ as
\[
R_A=\Tr(R_A)J+\sum_{j=1}^{N'}[S_j,T_j],
\]
where $J$ is a pseudodif\/ferential projection of rank 1 and $S_j$, $T_j$ are smoothing operators. Summing up, the expression for $A$ becomes
\begin{gather}
A=\sum\limits_{k=1}^{N}[\alpha_k,B_k]+\Tr(R_A)J+\sum_{j=1}^{N'}[S_j,T_j]. \label{commutators}
\end{gather}
Applying the canonical trace $\TR$ to both sides of this expression, since $\TR$ vanishes on commutators of operators in $\Cl_{\odd}(M)$, we infer that
\[
\TR(A)=\Tr(R_A).
\]
Thus \eqref{commutators} reads
\begin{gather}
A=\sum\limits_{k=1}^{N}[\alpha_k,B_k]+\TR(A)J+\sum_{j=1}^{N'}[S_j,T_j]. \label{commutators2}
\end{gather}
If $\tau$ is a trace on $\Cl_{\odd}(M)$, applying $\tau$ to both sides of~(\ref{commutators2}) we reach the conclusion of the theorem.
\end{proof}
\begin{thebibliography}{99}
\bibitem[3]{ducourtiouxthese}
Ducourtioux C., Weighted traces on pseudodif\/ferential operators and associated
  determinants, Ph.D. thesis, Universit\'e Blaise Pascal, Clermont-Ferrand,
  2001.

\bibitem[4]{fedosov}
Fedosov B.V., Golse F., Leichtnam E., Schrohe E., The noncommutative residue
  for manifolds with boundary, \href{http://dx.doi.org/10.1006/jfan.1996.0142}{\textit{J.~Funct. Anal.}} \textbf{142} (1996),
  1--31.

\bibitem[5]{grubb}
Grubb G., A resolvent approach to traces and zeta {L}aurent expansions, in
  Spectral Geometry of Manifolds with Boundary and Decomposition of Manifolds,
  \textit{Contemp. Math.}, Vol.~366, Amer. Math. Soc., Providence, RI, 2005,
  67--93, \href{http://arxiv.org/abs/math.AP/0311081}{math.AP/0311081}.

\bibitem[6]{guillemin}
Guillemin V., A new proof of {W}eyl's formula on the asymptotic distribution of
  eigenvalues, \href{http://dx.doi.org/10.1016/0001-8708(85)90018-0}{\textit{Adv. Math.}} \textbf{55} (1985), 131--160.

\bibitem[7]{guilleminfourier}
Guillemin V., Residue traces for certain algebras of {F}ourier integral
  operators, \href{http://dx.doi.org/10.1006/jfan.1993.1096}{\textit{J.~Funct. Anal.}} \textbf{115} (1993), 391--417.

\bibitem[8]{kontsevichdeterminant}
Kontsevich M., Vishik S., Determinants of elliptic pseudodif\/ferential
  operators, \href{http://arxiv.org/abs/hep-th/9404046}{hep-th/9404046}.

\bibitem[10]{lesch}
Lesch M., On the noncommutative residue for pseudodif\/ferential operators with
  log-polyhomogeneous symbols, \href{http://dx.doi.org/10.1023/A:1006504318696}{\textit{Ann. Global Anal. Geom.}} \textbf{17}
  (1999), 151--187, \href{http://arxiv.org/abs/dg-ga/9708010}{dg-ga/9708010}.

\bibitem[11]{leschtraces}
Lesch M., Pseudodif\/ferential operators and regularized traces, in Motives,
  Quantum Field Theory, and Pseudodif\/ferential Operators, \textit{Clay Math.
  Proc.}, Vol.~12, Amer. Math. Soc., Providence, RI, 2010, 37--72,
  \href{http://arxiv.org/abs/0901.1689}{arXiv:0901.1689}.

\bibitem[14]{maniccia}
Maniccia L., Schrohe E., Seiler J., Uniqueness of the {K}ontsevich--{V}ishik
  trace, \href{http://dx.doi.org/10.1090/S0002-9939-07-09168-X}{\textit{Proc. Amer. Math. Soc.}} \textbf{136} (2008), 747--752,
  \href{http://arxiv.org/abs/math.FA/0702250}{math.FA/0702250}.

\bibitem[18]{okikiolu}
Okikiolu K., The {C}ampbell--{H}ausdorf\/f theorem for elliptic operators and a
  related trace formula, \href{http://dx.doi.org/10.1215/S0012-7094-95-07918-6}{\textit{Duke Math.~J.}} \textbf{79} (1995), 687--722.

\bibitem[20]{ouedraogo}
Ouedraogo M.F., A symmetrized canonical determinant on odd-class
  pseudodif\/ferential operators, in Geo\-metric and Topological Methods for
  Quantum Field Theory, \href{http://dx.doi.org/10.1017/CBO9780511712135.012}{Cambridge Univ. Press}, Cambridge, 2010, 381--393.

\bibitem[21]{ouedraogothese}
Ouedraogo M.F., Extension of the canonical trace and associated determinants,
  Ph.D. thesis, Universit\'e Blaise Pascal, Clermont-Ferrand, 2009.


\bibitem[22]{sylvie}
Paycha S., The noncommutative residue and canonical trace in the light of
  {S}tokes' and continuity properties, \href{http://arxiv.org/abs/0706.2552}{arXiv:0706.2552}.

\bibitem[24]{scott}
Paycha S., Scott S., A {L}aurent expansion for regularized integrals of
  holomorphic symbols, \href{http://dx.doi.org/10.1007/s00039-007-0597-8}{\textit{Geom. Funct. Anal.}} \textbf{17} (2007),
  491--536, \href{http://arxiv.org/abs/math.AP/0506211}{math.AP/0506211}.

\bibitem[25]{PongeAsymmetry}
Ponge R., Spectral asymmetry, zeta functions, and the noncommutative residue,
  \href{http://dx.doi.org/10.1142/S0129167X06003825}{\textit{Internat.~J. Math.}} \textbf{17} (2006), 1065--1090,
  \href{http://arxiv.org/abs/math.DG/0310102}{math.DG/0310102}.

\bibitem[26]{ponge}
Ponge R., Traces on pseudodif\/ferential operators and sums of commutators,
  \href{http://dx.doi.org/10.1007/s11854-010-0001-8}{\textit{J.~Anal. Math.}} \textbf{110} (2010), 1--30, \href{http://arxiv.org/abs/0707.4265}{arXiv:0707.4265}.

\bibitem[28]{seeley}
Seeley R.T., Complex powers of an elliptic operator, in Singular {I}ntegrals
  ({P}roc. {S}ympos. {P}ure {M}ath., {C}hicago, {I}ll., 1966), Amer. Math.
  Soc., Providence, R.I., 1967, 288--307.

\bibitem[29]{Shubin}
Shubin M.A., Pseudodif\/ferential operators and spectral theory, 2nd ed.,
  Springer-Verlag, Berlin, 2001.

\bibitem[31]{wodzicki}
Wodzicki M., Noncommutative residue. {I}.~{F}undamentals, in {$K$}-Theory,
  Arithmetic and Geo\-met\-ry ({M}oscow, 1984--1986), \href{http://dx.doi.org/10.1007/BFb0078372}{\textit{Lecture Notes in
  Math.}}, Vol.~1289, Springer, Berlin, 1987, 320--399.

\bibitem[32]{WodzickiThesis}
Wodzicki M., Spectral asymmetry and noncommutative residue, Thesis, Steklov Institute, Soviet Academy of Sciences, Moscow, 1984.


\end{thebibliography}
\end{document}
